\section{Conclusion}
\label{sec:conclusion}

We cast harmful meme understanding as a structured,
evidence-grounded interpretation problem instead of only a binary
classification task.
The proposed framework jointly models harmfulness, target presence and
granularity, primary communicative intent, rhetorical tactics, and
multimodal relation while retaining an auditable record of the evidence
available to the prediction.

The framework decomposes this process into five stages:
internal multimodal evidence extraction, external knowledge acquisition,
knowledge verification, task-aware evidence fusion, and structured
interpretation with evidence attribution.
In particular, separating knowledge acquisition from verification prevents
retrieved or generated contextual information from being treated
automatically as trusted evidence, while provenance records preserve the
distinction among acquired, verified, rejected, and selected information.

We evaluate this formulation under a leakage-controlled cross-dataset
protocol in which models are developed only on HarMeme and evaluated on the
held-out Facebook Hateful Memes benchmark after model, retrieval, prompt,
and decoding decisions have been frozen.
Beyond the original harmfulness labels, schema-validated silver
annotations enable separate analysis of target, intent, rhetorical tactics,
and multimodal relation while making annotation coverage and provenance
explicit.

% TODO_RESULT:
% Replace this paragraph after the final audited result freeze.
% It should contain only 2--3 observations actually supported by the final
% tables, for example:
%
% Across the held-out target set, [Ours/result statement].
% The component analysis further shows that [retrieval/verifier/gate result].
% Structured evaluation indicates that [most/least transferable field].
%
% Do not introduce SOTA, statistical significance, or efficiency claims
% unless directly supported by the final audited experiments.

Taken together, this work provides a modular basis for studying not only
whether a meme is harmful, but also which structured interpretation the
model produces and what evidence was permitted to influence that decision.
The current results should nevertheless be interpreted in light of the
silver structured annotations, dataset-specific cultural coverage, and the
diagnostic rather than causal nature of the retained evidence and
rationales.
Future work can extend the protocol with independently adjudicated
structured annotations, broader cultural and linguistic domains, and
human-centered evaluation of evidence usefulness and rationale quality.
